\documentclass{article}
\usepackage{amsmath,amssymb}
\usepackage{xurl}
\usepackage[colorlinks=true,linkcolor=blue,citecolor=blue,urlcolor=blue]{hyperref}
% Shared commands for notes in docs/paper-gaps/.

\DeclareUnicodeCharacter{2080}{\ifmmode{}_{0}\else\textsubscript{0}\fi}
\DeclareUnicodeCharacter{2081}{\ifmmode{}_{1}\else\textsubscript{1}\fi}
\DeclareUnicodeCharacter{2082}{\ifmmode{}_{2}\else\textsubscript{2}\fi}
\DeclareUnicodeCharacter{2099}{\ifmmode{}_{n}\else\textsubscript{n}\fi}

\newcommand{\Lean}{\textsc{Lean}}
\ExplSyntaxOn
\NewDocumentCommand{\leanid}{m}{
  \ifmmode
    \textsf{\detokenize{#1}}
  \else
    \group_begin:
    \urlstyle{sf}
    \tl_set:Nn \l_tmpa_tl {#1}
    \tl_replace_all:Nnn \l_tmpa_tl {\_} {_}
    \exp_args:NV \path \l_tmpa_tl
    \group_end:
  \fi
}
\ExplSyntaxOff
\providecommand{\tr}{\operatorname{tr}}
\newcommand{\tnleanIssueBase}{https://github.com/LionSR/TNLean/issues/}
\newcommand{\tnleanPrBase}{https://github.com/LionSR/TNLean/pull/}
\newcommand{\tnleanDocBase}{https://sirui-lu.com/TNLean/docs/}
\newcommand{\ghissue}[1]{\href{\tnleanIssueBase#1}{GitHub issue \##1}}
\newcommand{\ghpr}[1]{\href{\tnleanPrBase#1}{PR \##1}}
\newcommand{\doclink}[1]{\href{\tnleanDocBase#1.html}{\path{#1}}}

% Dirac notation and the identity operator, as in blueprint/src/macros/common.tex.
% \providecommand keeps note-local definitions authoritative.
\usepackage{dsfont}
\providecommand{\ket}[1]{|#1\rangle}
\providecommand{\bra}[1]{\langle #1|}
\providecommand{\braket}[2]{\langle #1 | #2 \rangle}
\providecommand{\ketbra}[2]{|#1\rangle\!\langle #2|}
\providecommand{\Id}{\mathds{1}}
\providecommand{\one}{\mathds{1}}
\providecommand{\C}{\mathbb{C}}
\providecommand{\MN}[1]{M_{#1}(\C)}
\providecommand{\MD}{M_D(\C)}

% Verdict marker: \gapnote{<kind>}{<status>}. Kind is the policy
% classification (clarification, local-correction, scope-restriction,
% unfaithful, false-source, open-gap); status is open, wip (elimination
% underway), resolved, or historical. Machine-read by the site index;
% prints a verdict line.
\newcommand{\gapnote}[2]{\par\noindent\textbf{Verdict:} #1 (#2).\par}

\providecommand{\leanid}[1]{\textsf{\detokenize{#1}}}
\providecommand{\doclink}[1]{%
    \href{https://sirui-lu.com/TNLean/docs/#1.html}{\path{#1}}}
\providecommand{\gapnote}[2]{%
    \par\noindent\textbf{Verdict:} #1 (#2).\par}
\title{Conditional Final-Scale Information for the Physical Source}
\author{}
\date{October 14, 2026}
\begin{document}
\maketitle
\gapnote{scope-restriction}{open}

The concrete physical source now yields the final-scale information
estimate, conditional on its complete original hypotheses.  The proof
composes the same-witness information theorem with the uniform numerical
error bound.  It does not assume a target comparator, a terminal entropy
certificate, a selected density/energy conclusion, a mutual-information
bound or a safe-box improvement.

For $X=A\cap T$ and the same actual truncation set $S_0$, the conclusion is
\[
  2I_\Omega(X:S_0^c)\leq Cn^{62497/62500},\qquad
  2I_\Omega(X:S_0^c)\leq Cn^{1-1/100000}.
\]
The exact maximal exponent $62497/62500=0.999952$ has strict margin
$19/500000$ to $0.99999$.  The radius constant is fixed before the source
entropy coefficient $C_e$; the positive information coefficient and size
threshold are fixed after $C_e$ but before the finite physical system.
\footnote{Declaration:
    \leanid{TNLean.PEPS.AreaLaw.Scan.CollarScan.exists_sourceTarget_finalScale_mutualInformation}
    in \doclink{TNLean/PEPS/AreaLaw/Scan/SourceFinalInformation}.}

The same-source theorem is invoked once.  Its original and truncated
vectors, selected spectral window, label sequence and all actual selected
points remain unchanged.  Its replica limit was taken with the physical
size, weight and smoothing fraction fixed.  Only afterward is the final
size estimate applied at $W=n^{1/20000}$ and $\tau=n^{-1/4000}$.
The rounded-scale transport condition is proved eventually, not assumed.

The target dimension is derived from the same compact volume estimate:
$X\subseteq S_0$ and $|S_0|\leq(C_T+1)n^2$ imply
\[
  \log(e\dim\mathcal H_X)\leq1+((C_T+1)\log q)n^2.
\]
No far-region size is introduced.  The same trace-distance budget also
gives the target entropy transfer with error $n^{-247}$; the existing
shell transfer keeps its own $n^{-247}$ error.
\footnote{Declaration:
    \leanid{TNLean.PEPS.AreaLaw.Scan.log_exp_configuration_card_le_of_card}
    in \doclink{TNLean/PEPS/AreaLaw/Scan/PhysicalInformationDimension}.}

The following source obligations remain explicit:
\begin{itemize}
    \item The gapped Hamiltonian, physical scanner, induced domain graph,
          support anchors and fixed scanner coefficient.
    \item The literal radius, spacing, separation and band scales;
          $2r_0\leq D$, $1\leq D$ and $4D\leq m$.
    \item The nonempty template, quadratic volume, ambient depth,
          layer-cardinality, anchor-multiplicity and boundary-clearance bounds.
    \item The original-state target and every positive-prefix entropy
          input at $e_{\rm box}=2/1000000$, with the same fixed $C_e$.
\end{itemize}
This result does not construct those geometric or safe-box entropy
antecedents, close the remaining geometric information telescope, or prove
Proposition 9.6 or the full area law.  It is a conditional final-scale
physical information bound.  Hosted kernel, axiom and documentation checks
are required before acceptance; this scope note does not certify them.
\footnote{Source: Ref.~\cite{OpenAI2026AreaLaw}, \texttt{08-scanner.tex},
    lines 353--357, 399--533 and 839--941, revision
    \texttt{adc7f1241b42e322a6451854ab7e4b4c146bf78a}.}

\bibliographystyle{alpha}
\bibliography{references}
\end{document}
